% GENERATED FILE. Edit canonical lessons, then run npm run generate:books.
% canonical-source-hash: fnv1a64:e1dd1e34f6ab450b
% canonical-lessons: UR-C139-banana, UR-C139-torna, UR-C139-jalana, UR-C139-bhagna, UR-C139-bulana

\chapter{To make, To break, To burn, To run away, To call}
\label{ch:ur-to-make-to-break-to-burn-to-run-away-to-call}
\hlchaptermodality{139}

\begin{chapteropening}
\textbf{By the end of this chapter:} \emph{I can say to make, to break, to burn, to run away and to call.}

Complete the last lesson of chapter 139: I can say to make, to break, to burn, to run away and to call.
\end{chapteropening}

\section[\texorpdfstring{\ur{بنانا} (banānā)}{بنانا (banānā)}]{\ur{بنانا} (banānā) — to make}

\label{lesson:UR-C139-banana}



\begin{quote}
Before the new one: say the Urdu for to play, then the Urdu for to touch.
\end{quote}

\subsection*{You'll want to know: \ur{بنانا}}
\textbf{\ur{بنانا}} — \emph{banānā} — \textquotedblleft{}to make\textquotedblright{}.

\begin{cousinweb}
\textbf{\ur{بنانا}} (\emph{banānā}) is the causing form of \emph{bannā}, \textquotedblleft{}to be made\textquotedblright{}, which CDIAL lists under Sanskrit \emph{vanati}, \textquotedblleft{}desires, gains, makes ready\textquotedblright{} (CDIAL 11260).
\end{cousinweb}

\begin{grammarlens}[title={What you've built}]
One more word for this chapter.
\end{grammarlens}

\subsection*{Your turn}
\begin{itemize}
\raggedright
  \item \emph{Say it:} \emph{banānā}
  \item \emph{Say it:} \emph{banānā}, once more
  \item \emph{Say it:} say \emph{banānā}
  \item \emph{Recall:} say the Urdu for to play, then the Urdu for to touch, then say \emph{banānā} again
\end{itemize}

\begin{tcolorbox}[breakable,colback=teal!4,colframe=teal!35!black,title={Before you move on}]
What does \ur{بنانا} mean? (\textquotedblleft{}To make\textquotedblright{}.) Say it once more.
\end{tcolorbox}
\section[\texorpdfstring{\ur{توڑنا} (toṛnā)}{توڑنا (toṛnā)}]{\ur{توڑنا} (toṛnā) — to break}

\label{lesson:UR-C139-torna}



\begin{quote}
Before the new one: say the Urdu for to touch, then the Urdu for to make.
\end{quote}

\subsection*{You'll want to know: \ur{توڑنا}}
\textbf{\ur{توڑنا}} — \emph{toṛnā} — \textquotedblleft{}to break\textquotedblright{}.

\begin{cousinweb}
\textbf{\ur{توڑنا}} (\emph{toṛnā}) is a tadbhava, a word inherited from Sanskrit and reshaped by sound change. It comes from \emph{troṭayati}, \textquotedblleft{}breaks, tears\textquotedblright{} (CDIAL 6079).
\end{cousinweb}

\begin{grammarlens}[title={What you've built}]
One more word for this chapter.
\end{grammarlens}

\subsection*{Your turn}
\begin{itemize}
\raggedright
  \item \emph{Say it:} \emph{toṛnā}
  \item \emph{Say it:} \emph{toṛnā}, once more
  \item \emph{Say it:} say \emph{toṛnā}
  \item \emph{Recall:} say the Urdu for to touch, then the Urdu for to make, then say \emph{toṛnā} again
\end{itemize}

\begin{tcolorbox}[breakable,colback=teal!4,colframe=teal!35!black,title={Before you move on}]
What does \ur{توڑنا} mean? (\textquotedblleft{}To break\textquotedblright{}.) Say it once more.
\end{tcolorbox}
\section[\texorpdfstring{\ur{جلانا} (jalānā)}{جلانا (jalānā)}]{\ur{جلانا} (jalānā) — to burn, to light}

\label{lesson:UR-C139-jalana}



\begin{quote}
Before the new one: say the Urdu for to make, then the Urdu for to break.
\end{quote}

\subsection*{You'll want to know: \ur{جلانا}}
\textbf{\ur{جلانا}} — \emph{jalānā} — \textquotedblleft{}to burn, to light\textquotedblright{}.

\begin{cousinweb}
\textbf{\ur{جلانا}} (\emph{jalānā}) is the causing form of \emph{jalnā}, \textquotedblleft{}to burn\textquotedblright{}: to light is to make something burn. CDIAL lists both under Sanskrit \emph{jvalati}, \textquotedblleft{}burns brightly\textquotedblright{} (CDIAL 5306).
\end{cousinweb}

\begin{grammarlens}[title={What you've built}]
One more word for this chapter.
\end{grammarlens}

\subsection*{Your turn}
\begin{itemize}
\raggedright
  \item \emph{Say it:} \emph{jalānā}
  \item \emph{Say it:} \emph{jalānā}, once more
  \item \emph{Say it:} say \emph{jalānā}
  \item \emph{Recall:} say the Urdu for to make, then the Urdu for to break, then say \emph{jalānā} again
\end{itemize}

\begin{tcolorbox}[breakable,colback=teal!4,colframe=teal!35!black,title={Before you move on}]
What does \ur{جلانا} mean? (\textquotedblleft{}To burn, to light\textquotedblright{}.) Say it once more.
\end{tcolorbox}
\section[\texorpdfstring{\ur{بھاگنا} (bhāgnā)}{بھاگنا (bhāgnā)}]{\ur{بھاگنا} (bhāgnā) — to run away}

\label{lesson:UR-C139-bhagna}



\begin{quote}
Before the new one: say the Urdu for to break, then the Urdu for to burn.
\end{quote}

\subsection*{You'll want to know: \ur{بھاگنا}}
\textbf{\ur{بھاگنا}} — \emph{bhāgnā} — \textquotedblleft{}to run away\textquotedblright{}.

\begin{cousinweb}
\textbf{\ur{بھاگنا}} (\emph{bhāgnā}) is a tadbhava, a word inherited from Sanskrit and reshaped by sound change. It comes from \emph{bhagna}, \textquotedblleft{}broken, routed\textquotedblright{} (CDIAL 9361). It is the past form of \emph{bhajyate}, \textquotedblleft{}is broken\textquotedblright{}. Bloch linked the sense \textquotedblleft{}flee\textquotedblright{} with a different old verb; CDIAL finds that unnecessary, since \textquotedblleft{}broken\textquotedblright{} could come to mean \textquotedblleft{}routed, fled\textquotedblright{}.
\end{cousinweb}

\begin{grammarlens}[title={What you've built}]
One more word for this chapter.
\end{grammarlens}

\subsection*{Your turn}
\begin{itemize}
\raggedright
  \item \emph{Say it:} \emph{bhāgnā}
  \item \emph{Say it:} \emph{bhāgnā}, once more
  \item \emph{Say it:} say \emph{bhāgnā}
  \item \emph{Recall:} say the Urdu for to break, then the Urdu for to burn, then say \emph{bhāgnā} again
\end{itemize}

\begin{tcolorbox}[breakable,colback=teal!4,colframe=teal!35!black,title={Before you move on}]
What does \ur{بھاگنا} mean? (\textquotedblleft{}To run away\textquotedblright{}.) Say it once more.
\end{tcolorbox}
\section[\texorpdfstring{\ur{بلانا} (bulānā)}{بلانا (bulānā)}]{\ur{بلانا} (bulānā) — to call}

\label{lesson:UR-C139-bulana}



\begin{quote}
Before the new one: say the Urdu for to burn, then the Urdu for to run away.
\end{quote}

\subsection*{You'll want to know: \ur{بلانا}}
\textbf{\ur{بلانا}} — \emph{bulānā} — \textquotedblleft{}to call\textquotedblright{}.

\begin{cousinweb}
\textbf{\ur{بلانا}} (\emph{bulānā}) is the causing form of \emph{bolnā}, \textquotedblleft{}to speak\textquotedblright{}: to call someone is to make them speak back. CDIAL lists both under a reconstructed \emph{*boll-}, \textquotedblleft{}speak\textquotedblright{} (CDIAL 9321).
\end{cousinweb}

\begin{grammarlens}[title={What you've built}]
One more word for this chapter.
\end{grammarlens}

\subsection*{Your turn}
\begin{itemize}
\raggedright
  \item \emph{Say it:} \emph{bulānā}
  \item \emph{Say it:} \emph{bulānā}, once more
  \item \emph{Say it:} say \emph{bulānā}
  \item \emph{Recall:} say the Urdu for to burn, then the Urdu for to run away, then say \emph{bulānā} again
\end{itemize}

\begin{tcolorbox}[breakable,colback=teal!4,colframe=teal!35!black,title={Before you move on}]
What does \ur{بلانا} mean? (\textquotedblleft{}To call\textquotedblright{}.) Say it once more.
\end{tcolorbox}
